\chapter{Finanças}

\section{Fama-MacBeth}

O procedimento de Fama-MacBeth estima regressões de corte transversal
período a período e depois tira a média dos coeficientes resultantes.
Produz erros padrão robustos à correlação de corte transversal.

\begin{lstlisting}[caption={Regressão Fama-MacBeth}]
let n = 1000
generate df mes   = ceil(_n / 100)
generate df ret   = rnormal(n)
generate df beta  = rnormal(n)
generate df tamanho = rnormal(n)

let m_fmb = fmb(ret ~ beta + tamanho, df, time=mes)
print(m_fmb)
\end{lstlisting}

O comando roda uma regressão de corte transversal por período e tira a
média das inclinações, reportando a média temporal e o erro padrão de
cada coeficiente.

\section{Sorts de portfólio}

`portsort` forma portfólios ordenando ativos em quantis de uma
característica. `doublesort` faz isso em duas dimensões.

\begin{lstlisting}[caption={Sort de portfólio simples}]
let m_ps = portsort(df, ret, tamanho, n=5)
print(m_ps)
\end{lstlisting}

\begin{lstlisting}[caption={Sort duplo}]
let m_ds = doublesort(df, ret, tamanho, bm, n1=5, n2=5)
print(m_ds)
\end{lstlisting}

A saída reporta o retorno médio de cada célula de portfólio. Sorts
duplos são úteis para isolar a interação entre duas características.

\section{Notas práticas}

\begin{itemize}
\item Fama-MacBeth é apropriado quando betas são estimados em uma
      primeira etapa e depois usados como regressores em uma segunda.
\item Sorts de portfólio são descritivos, não causais. Mostram padrões,
      não identificação.
\item Use `if=` para excluir observações extremas ou focar em um
      período amostral.
\end{itemize}
